% =====================================================================
%  POINT-BY-POINT RESPONSE TO REVIEWERS
%  Manuscript: Quantifying comprehensive marine vessel emissions using satellite data fusion
%  MS #: NCOMMS-26-022400            Journal: Nature Communications
%
%  Conventions (per editor request, comments reproduced verbatim):
%   - Reviewer comments: gray left-barred boxes, reproduced verbatim.
%   - Author responses: blue "Author response:" label.
%   - Manuscript changes: green "Changes in manuscript:" label.
%   - Page/line references point to the track-changes (latexdiff) PDF.
%
%  WORKFLOW NOTE: keep the assignment/status bookkeeping in the Google Sheets tracker
%  sheet (https://docs.google.com/spreadsheets/d/1hIXrMYa4PrrSaBtfAJ_X6QiJfIspgWiX2U39bFFANgQ/edit?gid=0#gid=0); 
% this .tex file is the deliverable.
%  Each \point has a % category / % assignee / % bundled-subparts note
%  to help you triage against the sheet.
% =====================================================================
\documentclass[11pt]{article}
\usepackage[utf8]{inputenc}
\usepackage[margin=1in]{geometry}
\usepackage{xcolor}
\usepackage{mdframed}
\usepackage{enumitem}
\usepackage{parskip}

\definecolor{refbg}{gray}{0.95}
\definecolor{refrule}{gray}{0.55}
\definecolor{respblue}{RGB}{20,80,160}
\definecolor{changegreen}{RGB}{20,120,60}

\usepackage{hyperref}
\hypersetup{colorlinks=true, linkcolor=respblue, urlcolor=respblue}

% Verbatim reviewer comment box (gray, left bar)
\newmdenv[
  backgroundcolor=refbg, linecolor=refrule, linewidth=1pt,
  topline=false, bottomline=false, rightline=false, leftline=true,
  innerleftmargin=10pt, innerrightmargin=10pt,
  innertopmargin=6pt, innerbottommargin=6pt,
  skipabove=8pt, skipbelow=6pt
]{refquote}

\newcommand{\reviewer}[1]{\bigskip\section*{#1}}
\newcommand{\point}[1]{\subsection*{\normalfont\bfseries #1}}
\newcommand{\response}{\par\medskip\noindent%
  \textbf{\textcolor{respblue}{Author response:}}\ }
\newcommand{\changes}{\par\medskip\noindent%
  \textbf{\textcolor{changegreen}{Changes in manuscript:}}\ }
\newcommand{\todo}[1]{\par\smallskip\noindent\textcolor{red}{\itshape[TODO: #1]}}

\setlength{\parskip}{6pt}

\begin{document}

\begin{center}
  {\Large\bfseries Point-by-point response to the reviewers: Revision 1}\\[4pt]
  {\large \emph{Quantifying comprehensive marine vessel emissions using satellite data fusion}\\[2pt]
  Manuscript reference: NCOMMS-26-022400
\end{center}

\bigskip
\noindent We greatly appreciate the opportunity to revise our manuscript, and sincerely appreciate the thoughtfulness of the editor and referees in providing constructive feedback. Below we reproduce every comment verbatim (gray boxes)
and give our response in blue, with the corresponding manuscript changes in
green. All line and page numbers refer to the tracked-changes version of the
manuscript.

By incorporating this excellent feedback, we believe our manuscript is significantly stronger as a result.

We are also excited to highlight to the editor and both reviewers that since we first submitted the manuscript for review in March 2026, we have made a number of significant improvements to both the AIS-based and S1-based emissions models. For the AIS-based model, we have improved our vessel characteristic inference algorithms across the board, and are now able to distinguish between chemical and oil tankers. For the S1-based models, we have expanded our set of included predictive model features. The improved models all now have higher predictive power and out-of-sample performance. We have also increased the time period covered by our dataset to make our results more timely and globally relevant, expanding coverage that previously ended in 2024 to now include 2025. We have updated all manuscript figures, statistics, tables, and accompanying text to reflect these improved models and expanded temporal coverage.

\reviewer{Response to the Editor}

Before turning to the point-by-point responses to the referee comments, we would first like to highlight that we have fully formatted the manuscript in accordance with the guidelines provided by \textit{Nature Communications}:

\point{E1 --- Manuscript section order}

\response We have reordered the manuscript to follow the required structure. 

\todo{Two items in this list still need author text, and are flagged in the .tex:
(i) the Acknowledgements section is empty;
(ii) the Author Contributions statement is still the CRediT-style list and needs rewriting as one paragraph of descriptive sentences, as the editor asks}

\point{E2 --- Supplementary Information as a separate PDF}
% category: formatting | assignee: Gavin
% Separate PDF; cover page with title + full author info; items list on page 2;
% items labeled Supplementary Figure/Table/Note/Discussion/Methods/References.
\response \todo{confirm SI restructured with cover page + items list}

\point{E3 --- Figure labeling}

\response All multi-panel figures have been relabeled using lower case `a', `b', etc, using the same type as the other text in the figure.  All figure lettering is now also lower-case type, with only the first letter capitalized.


\point{E4 --- Color-vision-deficiency-safe figure colours}

\response We audited every figure in the manuscript to ensure each one is color-vision-deficiency-safe. Each categorical palette was simulated for deuteranopia, protanopia and tritanopia using the Machado et al. (2009) severity-1.0 model, and every pair of colours within a palette was then measured in CIE $\Delta E_{2000}$. For a categorical key the number that matters is the smallest pairwise distance, since that is the pair a reader is most likely to confuse.

Most of the manuscript already met the standard, because the figures draw their colours from palettes built for this purpose rather than chosen by hand: the Okabe-Ito palette for the data-source series that carry Figures 1, 2 and S1; Paul Tol's \emph{Safe} palette for categorical keys; the viridis family (\texttt{magma}, \texttt{mako}) for the gridded emissions maps; and Crameri's \texttt{scico} palettes for the diverging change maps. These are perceptually uniform and colour-vision-deficiency-safe by construction, and wherever a figure needed arbitrary colours we have taken them from one of these packages rather than inventing them. The training/testing map already used the turquoise-and-red
pairing the editor's guidance suggests. The audit also confirmed what these choices are worth: the gridded maps vary monotonically in lightness, so they carry their information without relying on hue discrimination at all, and the diverging change map holds 20.4 $\Delta E$ between its arms under protanopia.


\point{E5 --- Data Availability section}

\response We have added a Data Availability section.

\point{E6 --- Code Availability section}

\response We have added a Code Availability section.

\point{E7 --- Source Data file}

\response All manuscript figures, tables, and summary statistics can be reproduced entirely by the data and code provided in the GitHub repo, Zenodo repo, and Code Oceans repo. We therefore did not produce a separate `Source Data file'.

\point{E8 --- Reporting checklist(s)}
% category: policy | assignee: Gavin
% Must be completed in Adobe Reader (not a browser) and uploaded as a
% Related Manuscript file.
\response \todo{complete and attach}

\point{E9 --- ORCID for corresponding author(s)}
% category: admin | assignee: Gavin
\response \todo{link ORCID(s) in the MTS before acceptance}

\point{E10 --- Author-list change form (if applicable)}

\response The author list has not changed since our original submission.

% =====================================================================
\reviewer{Reviewer \#2}
% =====================================================================

\begin{refquote}\itshape
This manuscript proposes a novel global framework for estimating marine vessel
emissions by AIS data with satellite-(SAR). The authors address the limitation
of AIS-based approaches of incomplete data by incorporating detections of
non-broadcasting ships. The authors recommend a fusion approach to obtain
revised estimates of global maritime GHG emissions. Overall, the manuscript
tackles an important and timely problem with a technically sophisticated and
data-intensive approach. The contribution is significant both methodologically
and for policy applications. The manuscript is well-written, the results are
relevant for both scientific and policy communities. These issues need to be
addressed to improve the quality of the paper:
\end{refquote}
\response We thank the reviewer for the positive assessment and address each
point below.

\point{R2.1}

\begin{refquote}\itshape
The methodology relies on a key assumption that the relationship between
SAR-detected vessel density and emissions is similar for AIS-broadcasting and
non-broadcasting vessels (which may differ due to size, fuel types and
operational behavior). The authors could strengthen the justification by
conducting sensitivity analyses of different condition, and show potential bias
direction and magnitude more explicitly.
\end{refquote}
\response We thank the reviewer for bringing up this important point to clarify. We believe our model framework already addresses the majority of the reviewer's concern, and we appreciate the opportunity to clarify this for the reviewer and for future readers of the manuscript. For vessel sizes, we already bin the vessels and detections into fine length groups and drop small S1-unmatched detections below the smallest observed AIS-broadcasting vessel length, so we are not imposing this assumption to vessels of very different lengths. For operational behavior, our detections are already classified as fishing and non-fishing vessels, which is major factor determining operational behavior and emissions factors. For vessels of the same size and type, the fact that we are using Sentinel-1 imagery further helps reduce the variability inside each group, because SAR has limited ability of detecting wooden and fiberglass vessels. This means we are not applying the emission behavior of a small industrial fishing boat observed in the AIS data to a wooden small-scale fishing boat or fiberglass recreational fishing boat of similar size.

In addition, S1 vessel detections unmatched with AIS could correspond to \textit{either} vessels that do not carry an AIS transponder at all; \textit{or} to vessels that carry AIS transponders but temporarily non-broadcasting or whose signal was not received. Therefore, AIS-broadcasting vessels and unmatched vessel detections are not necessarily two distinct fleets, but can be two transient states. In these cases, the `relationship between SAR-detected vessel density and emissions' is exactly the same between the two transient states, because the two states represent the exact same set of vessels.

Finally, because our S1-unmatched model estimates emissions by space (1x1 degree pixels) and time (month) using AIS-broadcasting emissions data from the same space and time, the assumption is only applied to vessels operating in the same space and time. If vessels operating concurrently have similar characteristics, the assumption will hold. We do not use information from AIS-broadcasting vessels to infer S1-unmatched emissions in other places or times.

Nevertheless, we do still acknowledge there may be situations where the assumption breaks down.  This may be particularly the case when the emissions factors of non-broadcasting vessels are significantly different from the average emissions factors that form the basis of the 2020 Fourth IMO GHG report. For non-broadcasting vessels with newer technology that is significantly more efficient and less carbon intensive, this could bias our S1-unmatched emissions estimates too high. It is highly likely, however, that newer vessels with this improved technology would carry AIS transponders. For older non-broadcasting vessels that have outdated and more carbon intensive technology, this could bias our S1-unmatched emissions estimates too low. 

\changes We explicitly address this assumption, and why we believe it is largely valid, as its own paragraph in the discussion.

\begin{quote}
\textcolor{red}{Lines 677-695}: Several methodological limitations remain. One key assumption is that S1-unmatched emissions can be inferred from the relationship between S1 detections and AIS-based emissions for matched vessels. We believe this is a reasonable assumption for several reasons.  First, many unmatched S1 detections likely correspond to AIS-broadcasting vessels that are not observed in the AIS data because of transmission gaps or poor reception quality. In these cases, the relationship between S1 vessel detections and emissions would hold perfectly since these unmatched vessels are in fact just AIS-broadcasting vessels that happen to not be observed by AIS at a given point in time. Second, our model accounts for vessel length, a proxy for main engine power that is one of the primary determinants of emissions in AIS-based models and is applicable to virtually all vessels powered by internal combustion engines. Third, since fishing and non-fishing vessels are known to have different emissions factors, we account for this by differentiating detections between fishing and non-fishing vessels. Finally, our model incorporates spatiotemporal features that allow the relationship between detections and emissions to vary across space and time. Consequently, if broadcasting and non-broadcasting vessels operating in a given place and time have similar characteristics, the relationship between detections and emissions is expected to be similar as well. 
\end{quote}

Acknowledging there may be situations where the assumption breaks down, we also include the following discussion paragraph, where we discuss this possibility and how it would bias our estimates:

\begin{quote}
\textcolor{red}{Lines 696-706}: Another limitation is that the relationship between detections and emissions may fail to hold for non-broadcasting vessels if their emissions factors differ substantially from the IMO values we use to convert energy use into emissions. This could arise, for example, if non-broadcasting vessels disproportionately use newer and more efficient engine or fuel technologies that were not widely represented when the IMO emissions factors were developed in 2020; this would have the impact of biasing our S1-unmatched emissions estimates too high. Conversely, if non-broadcasting vessels instead use severely outdated and inefficient engine technologies, this this would have the impact of biasing our S1-unmatched emissions estimates too low. This highlights one of the primary limitations of the AIS model our methodology builds upon. Since information on the fuel type used by individual vessels is generally not publicly available, we rely on weighted average emissions factors from the IMO. Vessel-level emissions estimates based on average emission factors may be biased if a particular vessel's engine or fuel type is not well represented in the global fleet.
\end{quote}



\point{R2.2}

\begin{refquote}\itshape
The study relies on average emission factors from IMO guidelines. While standard
practice, this approach does not capture heterogeneity in fuel types and engine
technologies. Please clarify the implications of this assumption and discuss how
future work could incorporate vessel-specific fuel and engine data.
\end{refquote}
\response
Thank you for this clarifying comment. Since the initial submission, we have developed new machine learning models to predict engine type (oil, gas, or turbine), construction year, and engine speed (RPM) for each AIS-broadcasting vessel. This allows us to now more precisely assign emissions factors, including the ability to assign factors for different IMO NOx tiers and engine speeds. When a vessel's engine type cannot be resolved to a single emissions factor, we now use the predicted engine type to narrow the applicable emission factor columns and average across them. Although we attempted to develop a machine learning model to predict fuel type as well, we were unable to achieve reasonable predictive performance using the proprietary vessel registry we had access to. We now also discuss in the text how future research should be targeted at refining vessel-specific fuel and engine data.

\changes

We’ve added this new information in Section 4.1.2 in the following paragraphs:

\begin{quote}
\textcolor{red}{Lines 1338-1350}  Fuel consumption, for each of the main engine, auxiliary engine, and boiler engine, are next calculated as energy use multiplied by specific fuel consumption (SFC). The SFC for each of the main engine, auxiliary engine, and boiler engine are based on engine type and fuel type. Our vessel characteristics database includes broad categories of engine type (oil, gas, or turbine), as well as engine operating speed which allows us to differentiate between high speed diesel, medium speed diesel, and low speed diesel. Since our vessel characteristics database does not specify fuel type for each vessel, we use the average values from across fuel types for each type of engine. For the main engine, SFC is further modulated based on engine load: engines operating at low loads below 20\% operate inefficiently, and therefore consume more fuel for the same amount of energy output. For the two pollutants whose emissions factors depend on fuel use (CO$_{2}$ and SO$_{x}$), these emissions are therefore directly adjusted during inefficient low load operation.
\end{quote}

\begin{quote}
 \textcolor{red}{Lines 1370-1375} Again, since we have information for each vessel on engine type but not on fuel type, we take average values for the emissions factors for each engine type from across the appropriate fuel type options. NOx emissions factors additionally depend on the vessel’s IMO engine Tier, which is determined by year of build and engine RPM (both of which are included in our database of estimated vessel characteristics).
 \end{quote}

We have additionally added a description of the engine type, construction year, and engine speed machine learning models and their performance and interpretation to the methods section (\textcolor{red}{Lines 1135-1246: 4.1.6 Other vessel characteristics models'}). 

We finally added text in the Discussion noting the current inability to resolve fuel type and how additional training data could address this in future work:

\begin{quote}
 \textcolor{red}{Lines 743-755} This highlights one of the primary limitations of the AIS model our methodology builds upon. Since information on the fuel type used by individual vessels is generally not publicly available, we rely on weighted average emissions factors from the IMO. Vessel-level emissions estimates based on average emission factors may be biased if a particular vessel's engine or fuel type is not well represented in the global fleet. This could be an important consideration when using the model to explore policies aimed at technological engine changes and fuel switching. However, when this type of vessel-level information becomes publicly available, our modeling framework could easily be adapted to incorporate it. Despite the fact that we do not currently explicitly account for vessel-level fuel type, our vessel-level emissions model still achieves high performance when validated against two fully independent vessel-level datasets (Figs.  \textcolor{red}{6-7}).
\end{quote}

\point{R2.3}

\begin{refquote}\itshape
The manuscript uses multiple machine learning models (e.g., random forests) to
infer vessel characteristics and emissions. However, there is limited discussion
of model interpretability (e.g. feature importance) and potential overfitting or
bias propagation. The author could provide additional transparency, such as:
summary of key predictive features, robustness checks across model
specifications and brief discussion of interpretability limitations.
\end{refquote}

\response Thank you for the feedback. We agree that our initial draft lacked sufficient detail on the models. For each of the characteristic models, we now cover: (i) methodology, model application and training data, (ii) performance metrics, variable importance, and interpretability, and (iii) overfitting controls, known weaknesses, and potential biases. We believe that adding this level of detail is critical given the importance of vessel characteristics in AIS-based emissions models, and the novelty and scientific contribution of our methods for vessel characteristics inference.

\changes We added several dedicated section in the methods that includes a summary table of performance metrics for each model, alongside detailed breakdowns for each model with the information above:
\textcolor{red}{
\begin{itemize}
\item Lines 908-995: `4.1.2 Vessel class and characteristics models'
\item Lines 998-1055: `4.1.3 Primary GFW vessel class model'
\item Lines 1059-1096 4.1.4 Cargo sub-type model
\item Lines 1098-1132 4.1.5 Tanker sub-type model
\item Lines 1135-1246: 4.1.6 Other vessel characteristics models'
\end{itemize}}

We also made minor adjustments to the main text to point readers towards the newly written sections.


\point{R2.4}

\begin{refquote}\itshape
The authors can provide more precise the definition of ``non-broadcasting
vessels''.
\end{refquote}
\response Thank you for this suggestion. In reading through the suggestions from both reviewers, it has become clear to us that our terminology for specifying the two components of our emissions data fusion approach was confusing and imprecise. We have fully revised the text, and all figures and tables, to use a more clear set of terminologies: we now label emissions as either coming from either: 1) `AIS-broadcasting' vessels, or 2) `S1-unmatched vessels. Our use of the term 
`non-broadcasting vessels' was not clear or technically correct; this is because an S1 vessel detection that could not be matched to an AIS-broadcasting vessels  may correspond to \textit{either} a vessel not equipped with AIS, \textit{or} to a vessel equipped with AIS but experiencing a transmission gap.

\changes In the first paragraph of the results where we first introduce our results, we now explicitly state:

\begin{quote}
\textcolor{red}{Lines 209-215}: Throughout, we label the two components of this fusion explicitly: `AIS-based' emissions are those estimated from AIS-broadcasting vessels whose signal was received, and `S1-unmatched' emissions are those estimated from vessels detected by S1 that could not be matched to any AIS broadcast --- either because the vessel does not broadcast AIS at all, or because it does but its signal was not received. Neither term describes a vessel type or fleet: they name the observational basis of the estimate, and their sum is reported as `AIS-based + S1-unmatched'
\end{quote}

All figures and tables now also use `AIS-broadcasting' instead of `AIS'; and `S1-unmatched' instead of `non-broadcasting'.

We further clarify this in the methods section, `Sentinel-1 SAR data description':

\begin{quote}
\textcolor{red}{Lines 1606-1621}: By spatiotemporally matching AIS-broadcasting vessels to S1 detections, the dataset also categorizes each vessel detection as either being matched to an AIS-broadcasting vessel (`matched'), or being unmatched to an  an AIS-broadcasting vessel (`unmatched'). Importantly, although `unmatched' detections may often correspond to vessels that does not use AIS at all, it does not necessarily mean the vessel does not carry AIS - it could still carry AIS, but its transmission may be experiencing a gap for a variety of transmission or reception reasons. The detection-AIS matching utilizes a probabilistic framework that produces matching scores representing the likelihood of detection-AIS pairs being correct, and global threshold of the matching score. The matching considers AIS positions before and after the detection timestamp up to 12 hours away, and larger AIS position gaps typically result in lower matching scores with candidate detections. Therefore, an unmatched detection in this study may correspond to either a vessel not equipped with AIS or a vessel equipped with AIS but experiencing a transmission gap. The transmission gap can be intentional or a result of poor AIS reception
\end{quote}

We even go one step further. Since an `unmatched' detection could sometimes correspond to an AIS-broadcasting vessel that is simply in a transmission gap, we acknowledge that there is a potential for double-counting emissions from these vessels in both our AIS-broadcasting emissions estimate, and in our S1-unmatched estimate. We therefore have added a new analysis to estimate the potential magnitude of this double-counting and upward bias (it could be upwards of 2.8\% of our total fused estimated emissions). We describe this analysis in the methods, along with a new figure (Fig. \textcolor{red}{11}) and two new tables (Tables \textcolor{red}{6-7}):

\begin{quote}
\textcolor{red}{Lines 1621-1648}: A broadcasting vessel that goes unmatched in this way therefore erroneously has its emissions counted in both our `AIS-based' and `S1-unmatched' estimates, inflating our fused total. To estimate the extent of this double counting and its effect on our fused emissions estimates, we first use a hold-out experiment in which S1 detections with high-confidence AIS matches were re-matched using our matching algorithm but with AIS positions withheld within a window of the detection timestamp. This gives the probability that a detection of a broadcasting vessel is missed by the matching as a function of the AIS transmission gap; this rises gradually from 6.8\% at a 2-hour gap to 27.4\% at 8 hours and 45.4\% at 18 hours (Fig. \textcolor{red}{11a}). Second, because each AIS message represents the interval since the vessel's previous message, we know how much of our AIS-based CO$_{2}$ is emitted during gaps of each length (Fig. \textcolor{red}{6}, Fig. \textcolor{red}{11b}). Multiplying the two gives the share of AIS-based emissions emitted during gaps in which a detection of the vessel would have been missed, which we call the share at risk of double counting: 5.1\% (Fig. \textcolor{red}{11c-d}, Table \textcolor{red}{6}). Because S1 images a vessel at an arbitrary moment, this share is also the probability that any given detection of a broadcasting vessel is missed. Our 17.0 million matched detections are therefore the 94.9\% that were not missed, and the remaining 0.92 million detections of broadcasting vessels sit among our 14.9 million unmatched detections, where they make up 6.2\% of the total. Our S1-unmatched emissions are estimated from the density of unmatched detections, so a missed broadcasting vessel contributes to them like any other unmatched detection, and about 6.2\% of S1-unmatched emissions (171 million tonnes over 2017 to 2025, or 1.5\% of our combined total) therefore duplicate emissions already included in our AIS-based estimate. Finally, the missed detections are also absent from the matched detections used to train the emissions regression model, which raises the emissions it assigns to each detection by 5.4\%.  Upwards of 11.9\% of our S1-unmatched emissions may therefore be due to double-counting from missed matches during AIS transmission gaps; this translates to 2.8\% of our total fused emissions (Table \textcolor{red}{7}). 
\end{quote}

\point{R2.5}

\begin{refquote}\itshape
It would be useful to provide a simplified schematic or step-by-step explanation
and clearer separation between AIS-based and SAR-based components
\end{refquote}
\response We agree that the original schematic was difficult to follow. We have completely overhauled the schematic to more clearly step through the methods, and clearly separate the AIS-based and SAR-based components.
\changes The completely revised schematic can now be found in Fig. \textcolor{red}{5}.

\point{R2.6}
% category: figure | assignee: Gavin
% New figure: AIS-only vs fused estimates. Cross-ref R3.11 (spatial S1 map).
\begin{refquote}\itshape
A plot highlighting the differences of AIS-only vs fused estimates can improve
the quality of the paper.
\end{refquote}
\response Thank you for this comment and the opportuntity to improve our figures. In terms of annual temporal differences, Figure 1 panels A and B both show annual AIS-only emissions estimates (labeled `AIS-based'), estimates only from unmatched S1 detections (labeled `S1-unmatched'), and total fused estimates (labeled `AIS-based + S1-unmatched'). This figure shows how the data fusion approach increases the estimated magnitude of emissions relative to the AIS-only approach (panel A), while also dampening the increasing trend compared to the AIS-only approach (panel B). Similarly, the newly revised Figure 2 panel shows monthly AIS-only emissions estimates (labeled `AIS-based'), estimates only from unmatched S1 detections (labeled `S1-unmatched'), and total fused estimates (the stack of the two time series).
In terms of spatial differences, your comment also inspired us to revise Figure 2 to include a new Panel C that has spatial maps of AIS-only emissions estimates (top map, labeled `AIS-based'), estimates only from unmatched S1 detections (middle map, labeled `S1-unmatched'), and total fused estimates (bottom map, labeled `AIS-based + S1-unmatched'). We also have fully revised a companion figure to Figure 2, Figure S1, which is similar to Figure 2, but disaggregates temporal and spatial AIS-based, S1-unmatched, and fused emissions activity between non-fishing and fishing vessels. Taken together, Figures 1, 2, and S1 now clearly highlight both the temporal and spatial differences of AIS-only versus fused emissions estimates.

\changes Figures 1, 2, and \textcolor{red}{S1} have been revised.

% =====================================================================
\reviewer{Reviewer \#3}
% =====================================================================

\begin{refquote}\itshape
Overall, the manuscript is logically structured. However, the depth of the
results analysis and discussion is clearly insufficient. Consequently, although
the study adopts a new method, it does not appear robust, and the actual
contribution of this work to the field of ship emission inventories remains
unclear. My main concern is that, in order to highlight the novelty of the
proposed method, the authors have not conducted an objective literature review
and comparison with AIS-based methods. This leads to several potentially serious
misinterpretations for the research community, such as the estimated proportion
of ``dark ships'', which is not presented in an objective manner.
\end{refquote}
\response We thank the reviewer for these substantive comments. We address the overarching concern --- a fair, quantitative comparison with existing AIS-based inventories --- centrally in our responses to R3.1, R3.4 and R3.7. This comparison to other AIS-based inventories is now a central feature of our results, including its own paragraph and a new figure. It is also described in much more detail in a new extensive SI section. We believe this is an extremely important addition to the manuscript, and are grateful for your suggestions to incorporate this.

\point{R3.1}

\begin{refquote}\itshape
The manuscript develops a global ship emission inventory by integrating SAR and AIS data. The comparisons with a confidential test database containing ``measured fuel consumption'' for more than 15,000 ships and with the EU MRV annual CO\textsubscript{2} emissions reports are essentially validations at the individual ship level, and cannot by themselves support conclusions about global ship emissions. Given that many global ship emission inventory models based on AIS data already exist, such as SEIM, STEAM, MariTEAM, etc., the authors should provide a systematic comparison of their method, data (both quantity and quality), and results with existing methods, datasets, and inventories. This is necessary to demonstrate the added value and contribution of their new approach. In addition, the robustness of the method and the quality of the underlying data should be emphasized and discussed more clearly in the main text.
\end{refquote}

\response We thank the reviewer for this comment. Validation against measured fuel consumption and against EU MRV tests whether our model estimates emissions correctly at the individual ship level. However, it says nothing about whether the modeled fleet is complete, and therefore cannot on its own support a claim about global totals. We have retained those vessel-level validations, but we now frame them explicitly as tests of the emissions model rather than of the inventory, and we have added the fleet-completeness test the reviewer rightly asks for: a systematic comparison of our methods, our input data, and our results against the six most recent and representative global bottom-up, AIS-based inventories (STEAM, SEIM, MariTEAM, the OECD bottom-up model, the ICCT's SAVE model, and the model used for the Fourth IMO GHG Study).

We have structured this along exactly the three axes the reviewer names.

1. Method. Every AIS-based inventory is assembled from four components: the AIS record, a vessel-characteristics dataset, an engine model, and fuel consumption and emission factors. We compare ours against all six inventories on each. The comparison is transparent about both our limitations and the advantages of our broader scope. On the engine model, we use a load-based formulation, as do SEIM, SAVE, OECD and the Fourth IMO GHG Study, while STEAM and MariTEAM are resistance-based; we explain why the load-based choice is the appropriate one given that we model a fleet dominated by low information vessels, and we address the reported tendency of load-based models to yield higher fleet-level emissions than resistance-based. On fuel and policy treatment we are explicitly narrower than several of the others: we do not model fuel type at vessel level, we exclude LNG, and some specific emissions limits are assigned by engine Tier but not resolved spatially by ECA boundaries.

2. Data quantity and quality. On quantity, restricting the comparison to vessels explicitly modeled from observed AIS activity, we model 864,738 unique vessels, an average of 396,221 per year, which exceeds the AIS-modeled fleet of the Fourth IMO GHG Study, OECD, SAVE, MariTEAM and SEIM. Our fleet is comparable only to STEAM, though STEAM only reports fleet size for 2015. Our AIS record comprises approximately 125 billion positions over the study period after cleaning, which is only directly comparable with the yearly AIS messages used for STEAM, although STEAM reports message counts only for 2015. Part of the higher density in our later-period data comes from Dynamic AIS (D-AIS, or shipborne AIS), whereas other inventories only report the use of conventional AIS data. On quality, the most direct evidence is a like-for-like comparison of validation. The recent STEAM ambient-conditions update validates modeled fuel consumption against 6,445 vessel-year records in EU MRV. We validate against 20,432 vessel-years of the same benchmark. Both achieving similar performances. We also compare how every other inventory has been validated, which is mostly against aggregate statistics (IEA, UNCTAD, EMEP, national Air Emission Accounts) or against other inventories. We validate our aggregate results against all six bottom-up inventories and two top-down inventories (see response to R3.7).

3. Results. We compare our estimated magnitudes and trends against all six inventories, and separately against the fuel-statistics-derived EDGAR and CEDS. This is the comparison the reviewer asks for in R3.7, and we describe it in full there. We also discuss the spatial and temporal resolution of other inventories and their availability relative to ours.

On the added value of the approach, the contribution of our inventory lies in specific features of its scope and methodology. Existing inventories largely rest on a registry-matched core of large vessels that accounts for most emissions. They differ in how much of the low-information tail beyond this core they include, how they characterize these vessels, and how they treat vessels absent from AIS. Our AIS-based fleet is defined by AIS observations rather than registry coverage, allowing us to model emissions across a larger number of observed vessels, whereas, with some exceptions, other inventories remain largely bounded by a registry-defined core. We further contribute by estimating emissions for vessels that appear in neither registries nor the AIS record, using Sentinel-1 SAR detections that cannot be matched to an AIS broadcast. To our knowledge, this is the first global shipping inventory to do so.

\changes We have added the following.

1. Supplementary Note 2, `Systematic comparison of our AIS-based emissions modeling methodology to other AIS-based inventories'. This is the principal new material responding to this comment. It opens by situating the two methodological families (top-down fuel-based and bottom-up activity-based) and identifying the six inventories compared, then proceeds through: scope of the comparison and modeling approaches; activity and vessel-characteristics data (AIS coverage, fleet size, message volume, temporal coverage, how each inventory sources and gap-fills vessel characteristics, and how each treats vessels absent from AIS); the engine and emissions model (load-based versus resistance-based, main engine load corrections, SFOC, auxiliary and boiler decomposition, fuel types, and ECA and NOx Tier treatment); and differences in validation.

2. Supplementary Data 1, 2 and 3 (CSV files), giving the full attribute-by-attribute comparison with one column per inventory, so that readers can audit the comparison rather than take our summary of it:
\begin{itemize}
\item Supplementary Data 1, \path{bottom-up_inventory_comparison_methods.csv} --- main engine power and engine types, fuel types, main engine correction factors, specific fuel oil consumption, low-load adjustment, auxiliary engine and boiler power, emission control areas, imputation for vessels without AIS data, and validation.
\item Supplementary Data 2, \path{bottom-up_inventory_comparison_inputs.csv} --- AIS receiver types, vessel-characteristics sources, spatial and temporal coverage, number of vessels, and number of AIS messages after processing.
\item Supplementary Data 3, \path{bottom-up_inventory_comparison_outputs.csv} --- pollutants and other outputs, spatial and temporal resolution, disaggregation dimensions, and data access.
\end{itemize}

3. Supplementary Note 3 and the accompanying new main-text figure (Fig. 4), Results paragraph and Discussion paragraph, which carry the comparison of results. These are described in our response to R3.7.

4. Robustness and data quality in the main text. Addressing the reviewer's closing point, the main text now incorporates this material rather than leaving it to the Supplementary Information. The Results contain a dedicated paragraph benchmarking our estimates against other inventories, while the Discussion contains a paragraph interpreting the comparison and clarifying what it does and does not establish. In addition to the validation section within Methods, the main text now also reports details on the data sources, data quality, and inference of vessel characteristics. We have also added a new Methods section quantifying potential double-counting between our AIS-based and S1-unmatched components, as well as a new Supplementary Note 1 on trends in AIS receiver types.


\point{R3.2}
% category: MAJOR / analysis+framing | assignee: Anna
% Reviewer challenges the 84% "no registry info" claim and the resulting
% "AIS systematically underestimates" framing. Likely needs both a data
% response AND softened claims. Cross-ref R3.4.
\begin{refquote}\itshape
The manuscript states that approximately 84\% of AIS vessels lack publicly
available registry information and therefore require machine-learning-based
inference of main engine power and design speed. However, ship emission
calculations generally do not rely on the vessel information fields embedded in
AIS messages; instead, dedicated ship databases are used. Ship registry
information can be obtained through multiple channels, and although the coverage
of different databases varies, it is unlikely to be as low as suggested by the
authors. The issue may be that the authors did not obtain or use such databases;
it is therefore not appropriate to generalize from this limitation and claim
that AIS-based approaches systematically underestimate emissions. In the
manuscript, roughly 10\% of vessels with IMO registration are directly matched
to registry data, and another $\sim$6\% are matched via public ship databases
(non-IMO registered). This seems to reflect limitations of the authors' specific
implementation rather than an inherent flaw of AIS-based methods.
\end{refquote}
\response The reviewer is right, and we are grateful for the push. Our original sentence did not describe what we had actually measured. The 84\% came from a single field recording whether a vessel appears in any registry we could match to; it was not a measurement of whether main engine power and design speed specifically had to be inferred. Presenting the former as though it were the latter, in the same sentence, was our error, and we have rewritten the passage.

The arithmetic has also moved since submission, because the fleet is now defined over 2017--2025 and the vessel-characteristics database has been rebuilt. Of the 775,358 AIS-broadcasting vessels in the current analysis, 93,379 (12.0\%) carry an entry in IMO's GISIS, a further 45,743 (5.9\%) appear in at least one other registry, and 139,122 (18\%) are matched in total. We should also correct one small misreading: the 10\% in the submitted text was the share of all AIS-broadcasting vessels carrying an IMO entry, not the share of IMO-registered vessels we were able to match. Among vessels that are IMO-registered, our match rate is essentially complete. We would also flag, since it works against us, that these counts pool the whole study period: because registered vessels broadcast persistently while unregistered vessels appear intermittently and may change MMSI, the pooled share is less favorable than any single year's. Within individual years the registry-matched share runs from 30\% of broadcasting vessels in 2017 to 17\% in 2025. We now report this in the Results.

On the substance of the concern, that we did not obtain or use dedicated ship databases: we did, and we accept that the submitted text did not make this visible, citing only a single reference for the whole enterprise. Vessel characteristics in this analysis are drawn from IMO's GISIS; from S\&P Global Market Intelligence's Ships Dataset; from roughly thirty national and regional registries, among them those of the United States, Canada, Australia, Russia, Norway, the Republic of Korea, Indonesia, Iceland, the Faroe Islands, Spain and Taiwan; from the EU fleet register and FAO records; and from the authorized-vessel lists of the regional fisheries management organizations (ICCAT, IOTC, WCPFC, IATTC, SPRFMO, CCSBT, NEAFC, NPFC and GFCM). Across our full vessel-characteristics database of 864,738 vessels, 107,078 carry a GISIS listing, 23,793 and 12,824 appear in the two United States registries, 12,380 in TMT records and 10,679 in the EU register, with the RFMO lists contributing between roughly 1,300 and 4,300 vessels each. We have added this enumeration to the Methods.

The reason coverage is nevertheless low is not database access but fleet definition, and this is the point we should have made at the outset. Commercial ship databases are built around the IMO-registered merchant fleet, and for that fleet their coverage is excellent, which is precisely why we match it essentially completely. Our fleet, however, is defined by AIS observation rather than by registry membership. The 636,236 vessels we cannot match are overwhelmingly outside the merchant fleet those databases describe: on our classification roughly 43\% are passenger craft, a deliberately broad class that includes yachts, small ferries and Ro-Pax, 24\% are fishing vessels, and 9\% are tugs. No commercial database offers systematic coverage of these vessels, and acquiring further databases would not materially change the figure. We would note that this is a statement about the composition of the AIS-broadcasting fleet, not a claim about the quality of any particular database.

This is also why the per-characteristic figures are lower than 84\%, not higher. The new Methods table reports availability for each characteristic the emissions model requires. Main engine power is available from a registry or a commercial database for 7.5\% of vessels, and design speed for 1.8\%; the remainder are inferred. We have replaced the single 84\% figure with this table, so that the coverage claim is now made per characteristic rather than through a proxy for it.

Finally, we accept the framing point and have acted on it. The claim that AIS-based approaches systematically underestimate emissions should not rest, and no longer rests, on our registry coverage. Our 21\% figure is the share of the fused total contributed by Sentinel-1 detections that could not be matched to any AIS broadcast; it concerns vessels absent from AIS, not vessels absent from registries, and does not depend on our registry coverage in any way. The rewritten Results passage no longer runs the two together. The separate and considerably weaker claim, that an inventory whose fleet is defined by registry membership cannot represent vessels outside those registries, is now stated explicitly as a matter of fleet scope. We have removed the implication that this is an inherent flaw of AIS-based methods rather than a consequence of how a given inventory defines its fleet. This connects to our response to R3.3 on the choice of AIS provider and pipeline.

\changes We have rewritten the registry-coverage passage in the Results, added the per-characteristic availability table to the Methods (``Availability of vessel characteristics used in emissions modeling, by source'') together with the enumeration of registry and database sources above, added the within-year coverage range, and labeled each of the three fleet counts used in the paper (775,358 vessels emitting over 2017--2025; 864,738 vessels modeled over 2015--2025 in the inventory comparison; 1.8M MMSIs scored by the vessel class model) so that they can no longer be read against one another. The revised Results passage reads:

\begin{quote}
Of the roughly 775,000 AIS-broadcasting vessels in our analysis (every vessel emitting at any point during our 2017--2025 study period), we matched 139,000 (18\%) to public registries or proprietary databases, the majority of which (93,400) were IMO-registered. This allowed us to leverage vessel characteristics provided for these vessels from these databases. However, the vast majority of AIS-broadcasting vessels were not present in these registries, meaning that we needed to use model-derived vessel characteristics for between 83\% and 92\% of vessels, depending on the characteristic. Because this count pools the full study period, and because registered vessels broadcast persistently while unregistered vessels appear intermittently, it is not the coverage of the fleet in any one year: the registry-matched share of broadcasting vessels falls from 30\% in 2017 to 17\% in 2025, as the AIS-observed fleet grows faster than the registered fleet.
\end{quote}
\todo{Fill in the compiled line numbers for the quoted Results passage, and insert the Methods table number (expected Table 3, after the message-hour-threshold and model-summary tables) once Overleaf has rebuilt.}


\point{R3.3}
% category: discussion/framing | assignee: David
% Acknowledge + discuss impact of AIS provider/pipeline choice (GFW-only).
\begin{refquote}\itshape
The authors use only the GFW AIS dataset. In practice, the choice of AIS data
provider and processing pipeline has a substantial impact on data quality and
coverage. Many of the claimed advantages of the proposed method appear to arise
because the authors have not fully explored the broader AIS-based methodological
framework and available AIS datasets. This point should be acknowledged and
discussed.
\end{refquote}

\response We have documented the AIS sourcing and processing in the Methods. Our original text described GFW's AIS inputs in a single sentence, which was not adequate given how much the reviewer rightly says rests on it. Section 4.1.1 (AIS data description) now names the provider of our terrestrial, satellite and dynamic AIS, states that the feed covers both Class A and Class B transponders, distinguishes position messages from static and voyage messages, and describes the removal of invalid, duplicated and physically implausible positions before any downstream processing. It then describes how the retained pings are grouped into segments, which determines the time interval each ping represents and therefore the activity over which emissions are estimated \cite{kroodsma2018tracking}. We have also included a Supplementary Note 1, where we decompose and discuss our emissions, message volume and vessel counts by receiver type across the study period.

We have also compared our AIS inputs against the other bottom-up inventories as far as their reporting allowed. Supplementary Note 2 and Supplementary Data 2 set our AIS record against all six on the attributes the others report.

\changes We have made the following additions.

1. An expanded description of the AIS data and processing in the Methods, Section 4.1.1 "AIS data description". 

2. Supplementary Note 1, `Trends in AIS receiver types', with three supporting supplementary figures.

3. A comparison of our AIS inputs against the other inventories, in
Supplementary Note 2 and Supplementary Data 2 (\path{bottom-up_inventory_comparison_inputs.csv}).

\point{R3.4}
% category: MAJOR / analysis+framing | assignee: Zihan and David
% Overstating SAR advantage / understating AIS cleaning+reconstruction.
% Strong push for objective data-based comparison. Ties to R3.1, R3.2.
\begin{refquote}\itshape
Most advanced ship emission models based on AIS undertake extensive data
cleaning and correction, including enriching AIS with external vessel databases,
correcting erroneous positions, and interpolating trajectories. As a result, the
amount of ``missing'' emissions is generally much lower than implied in the
manuscript, and many supposedly ``dark'' ships are actually vessels with
erroneous or incomplete AIS trajectories that can be partially reconstructed. The
manuscript appears to overstate the advantages of the proposed SAR-based method
while downplaying the capabilities of existing AIS-based approaches. The authors
are strongly encouraged to conduct an objective and detailed data-based
comparison with existing global ship emission inventories.
\end{refquote}

\response Estimating emission from only AIS misses contributions from both truly non-broadcasting vessels and broadcasting but signal-not-received vessels. In theory, if we can reconstruct all the AIS activities of the latter, we can improve the AIS-derived emissions and get closer to the ``truth''. In practice, however, no publicly reproducible pipeline can truly collect, clean and reconstruct all the AIS activities, especially from publicly available data. By using processed AIS data from different providers, one can get different results. Regardless of the choices of reconstruction methods and AIS data providers, there will be bias in the AIS-based emission estimation. This bias and the missing contribution from truly non-broadcasting fleets together form the gap between comprehensive marine vessel emissions and any AIS-derived emissions.

We argue that the major value of approach is to provide a flexible solution for this challenge - it tolerates incomplete AIS data. The contribution by the unmatched vessel detections captures the emissions by both truly non-broadcasting vessels and broadcasting vessels with missing AIS data. Our model framework bypasses the requirement for the ``perfect'' AIS-derived emissions and directly leads to comprehensive marine vessel emissions.

\changes In both Introduction and Discussion sections, we added that AIS-based emission inventories can mitigate AIS data issues through data cleaning, registry enrichment, and trajectory reconstruction, but the limitation of non-broadcasting vessels and AIS coverage remain. We added discussion that the major value of approach is to provide a flexible solution for this challenge - it tolerates incomplete AIS data.

\point{R3.5}
% category: analysis/uncertainty | assignee: Zihan and Anna
% Line 410: non-vessel AIS signals in unmatched data; quantify or acknowledge.
% Cross-ref R2.4 (definition of non-broadcasting vessels).
\begin{refquote}\itshape
Line 410: AIS data contain many non-vessel signals, such as those from fishing
gear, aids to navigation, lighthouses, and other stationary objects. Therefore,
not all unmatched AIS signals correspond to actual vessel transmissions. The
authors should discuss the possible proportion of non-vessel AIS signals within
the unmatched AIS data. If such information cannot be quantified, this should at
least be explicitly acknowledged as a source of uncertainty.
\end{refquote}

\response Thank you for raising this important point. Roughly 5\% of our AIS-enabled signals are attributable to buoys, FADs, navigation aids, offshore structures, fishing gear, and other objects that are not vessels. Our primary AIS vessel classification model labels these non-vessels as `gear', and we remove them completely from our analysis.  In addition to removing non-vessel signals from AIS data, we have also removed possible non-vessel objects from SAR vessel detections by excluding detections around known infrastructures and repeated detections at fixed locations. Fishing gears are too small to be detected from Sentinel-1 SAR imagery. We now describe this at various places in both the results and methods sections.

\changes In the main results text when describing AIS-broadcasting emissions by different vessels classes, we have added the following:

\begin{quote}
\textcolor{red}{Lines 411-416:} We note that roughly 5\% of AIS-enabled signals are actually attributable to AIS transponders that are not on vessels, but are rather on buoys, fish aggregating devices (FADs), navigation aids, offshore structures, fishing gear, and other objects. These non-vessel AIS signals are completely excluded from our analysis.
\end{quote}

In the new methods Section, \textcolor{red}{4.1.2 Vessel class and characteristics models}, where we describe our vessel type classification models, we also now address this:

\begin{quote}
\textcolor{red}{Lines 1034-1038:} Note that roughly 5\% of AIS signals are likely associated with non-vessel objects (such as buoys, offshore structures, FADs, etc). These objects are identified as their own vessel class `gear' by the primary GFW vessel class model, and are excluded from the emissions estimates in this paper.
\end{quote}

In Section \textcolor{red}{4.2.1 Sentinel-1 SAR data description}, where we describe our S1 vessel detection processing procedure, we say:

\begin{quote}
\textcolor{red}{Lines 1549-1556:} The data filtering involved in the GFW vessel detections include removing detections within a 1-km buffer from the shoreline to reduce false positives from land-based objects due to inaccurately defined and dynamic coastlines, repeated detections of potentially fixed objects such as offshore infrastructure, detections from icy areas, moving vehicles close to shore, and other radar ambiguities. Non-vessel fishing gears are too small to be detected from Sentinel-1 SAR imagery, and thus no special filtering is necessary to remove these from the dataset.
\end{quote}

\point{R3.6}
% category: MAJOR / analysis | assignee: Zihan
% S1 orbital/temporal sampling biases vs dynamic ship activity; monthly aggregates.
% Editor's flagged "robustness" concern.
\begin{refquote}\itshape
The orbital characteristics of Sentinel-1 and the sampling properties of the SAR
data require further discussion. SAR observations are highly dependent on the
spatial and temporal coverage of the satellite passes, whereas ship activities
are highly dynamic. The current study appears to use monthly aggregates, using
limited overpass times to infer density and represent overall activity. However,
ship movements do not naturally follow a monthly timescale. The authors should
assess and discuss the potential biases and uncertainties introduced by this
temporal sampling, particularly for reconstructing ship trajectories and
estimating emissions.
\end{refquote}
\response We thank the reviewer for highlighting the fundamental mismatch between the sparse temporal sampling of Sentinel-1 and the highly dynamic nature of vessel movements. If we were applying a naive linear scaling (e.g., multiplying instantaneous SAR counts by the hours in a month), this temporal sparsity would indeed introduce severe biases.

However, our methodology is specifically designed to overcome this orbital limitation through our multi-model framework described in section 4.2.2. By training the model to predict AIS-broadcasting emissions using matched detection density, the model intrinsically learns to account for Sentinel-1’s revisit frequency and diurnal sampling biases. Consequently, when predicting emissions for non-broadcasting vessels, the model does not assume a steady-state temporal density from a single snapshot; rather, it applies the empirical temporal scaling learned from the continuous behavior of the AIS-broadcasting fleet. This model design is also described in the Introduction section.

\changes At the beginning of section 4.2.2, we have added explicit description of the potential bias that the Sentinel-1 sampling could introduce, which is a reason for developing our model framework.

\point{R3.7}
% category: analysis | assignee: Pol
% EDGAR as sole benchmark insufficient; add shipping-focused inventories.
% Ties directly to R3.1 -- same comparison exercise, different benchmark.
\begin{refquote}\itshape
EDGAR 2025 is a multi-sector emission inventory and shows noticeable differences in total emissions of various pollutants compared with specialized ship emission models. Using EDGAR as the sole or primary benchmark is therefore not sufficient. The authors are advised to include comparisons with additional shipping-focused inventories and models.
\end{refquote}
\response We fully agree, and thank the reviewer for pushing us on this. Benchmarking only against a multi-sector top-down inventory was indeed insufficient, and it also left the comparison uninformative about the thing our paper actually claims. We have therefore replaced the single EDGAR benchmark with a systematic comparison against the six global bottom-up, AIS-based shipping inventories (STEAM, SEIM, MariTEAM, the OECD bottom-up model, the ICCT's SAVE model, and the model used for the Fourth IMO GHG Study) and we retain the top-down inventories as a separate, clearly labeled second comparison, to which we have added CEDS alongside EDGAR. This comparison has its own paragraph in the Results, a new main-text figure, its own paragraph in the Discussion, and two new extensive Supplementary Notes, supported by three new Supplementary Figures, three new Supplementary Tables, and three Supplementary Data files.

Two points of clarification on how we have structured this. First, given its length, all the details and extensive text on this topic are included in the Supplement. Supplementary Note 2 compares methods and input data against the six shipping-focused inventories (this is the comparison the reviewer asks for in R3.1 as well). Supplementary Note 3, which is the direct answer to this comment, compares our estimated magnitudes and trends against those same six inventories first, and against EDGAR and CEDS. Second, the main highlights are summarized in the main text in the Results and Discussion, with the relevant notes and supplementary materials referenced.

The comparison is also considerably more informative than the original EDGAR benchmark, and we believe it substantially strengthens the paper. Our AIS-based estimates agree closely with the shipping-focused inventories over the early years of our study period and diverge in later years, and the new analysis identifies why: global AIS-based emissions are likely understated whenever heavy reliance on registry-defined fleets leaves low-information vessels out of scope. The vessels in our analysis are not limited to only those that are registered, but rather includes all vessels observed by AIS; it thus expands with both improvements in AIS technology along with the increasing uptake of AIS among vessels, with our estimated emissions growth mostly carried by low-information vessels. This indicates that the other bottom-up inventories, and our own early years, likely miss part of the global fleet. We also now report explicitly that this cuts both ways: because our fleet is strongly influenced by AIS observation, part of our total measured growth reflects expanding AIS carriage rather than rising activity at sea, and we say so in the Results, the Discussion and Supplementary Note 3. 

Finally, on the reviewer's specific point about EDGAR: our fused total does exceed both EDGAR and CEDS, and we now explain the three mechanisms responsible rather than simply reporting the gap. Fishing largely falls outside the shipping sector in both; small coastal and inland vessels are only partly covered; and neither inventory can capture fuel that never passes through reported channels.
\changes This comment prompted the following additions:

1. A new Results paragraph and a new main-text figure. The Results now include a dedicated paragraph benchmarking our estimates against the shipping-focused inventories first and the multi-sector inventories second, supported by new Fig. The new figure plots our AIS-based and fused annual CO$_{2}$ series against all six bottom-up, AIS-based inventories (panel a) and against the EDGAR and CEDS shipping sectors (panel b), so that the shipping-focused benchmarks are now the primary comparison and the multi-sector ones are clearly separated as a secondary check.

2. A new Discussion paragraph interpreting the comparison, and stating plainly what it does and does not establish:

3. Two new Supplementary Notes. Supplementary Note 3, `Systematic comparison of our emissions magnitude and trend estimates to other inventories' is the full version of this comparison. Supplementary Note 2, `Systematic comparison of our AIS-based emissions modeling methodology to other AIS-based inventories', provides the methodological counterpart that makes the results comparison interpretable. We describe this Note more fully in our response to R3.1.

4. Three new Supplementary Figures supporting the comparison.
5. Three new Supplementary Tables tabulating values discussed in the comparison.
6. Three new Supplementary Data files giving the full attribute-by-attribute comparison of methods, inputs and outputs against all six inventories bottom-up, one column per inventory: \path{bottom-up_inventory_comparison_methods.csv} (Supplementary Data 1), \path{bottom-up_inventory_comparison_inputs.csv} (Supplementary Data 2), and \path{bottom-up_inventory_comparison_outputs.csv} (Supplementary Data 3).


\point{R3.8}
% category: writing/minor | assignee: Zihan
% Intro line 131: explain what S1 SAR is and why it complements non-broadcasting.
\begin{refquote}\itshape
Introduction, line 131: when Sentinel-1 (S1) Synthetic Aperture Radar (SAR) is first introduced, the authors should briefly explain what this dataset is and why it can be used to complement emissions from non-broadcasting vessels. This will help readers who are not familiar with SAR to understand the rationale of the approach.
\end{refquote}
\response We thank the reviewer for the suggestion. In the introduction when we first introduce S1 SAR, we added a brief description of the S1 imagery dataset, including the resolution, the all-weather imaging ability, and the revisit frequency.

\changes We added a the following description to the introduction where we first introduce S1 SAR:

\begin{quote}
\textcolor{red}{Lines 139-148:} Here we develop a method to estimate high-resolution marine vessel emissions globally from 2017 through 2025 by fusing AIS data from broadcasting vessels with Sentinel-1 (S1) Synthetic Aperture Radar (SAR) detections of vessels that are not broadcasting AIS. The 20m-resolution S1 SAR imagery allows us to detect and estimate the lengths of most vessels larger than 15 m regardless of whether they broadcast AIS, cloud cover, weather conditions, or solar illumination. The two S1 satellites active during our study period (full period for S1A, 2017-2021 for S1B) operate in a polar, sun-synchronous orbit, providing a global revisit frequency of 6 to 12 days. S1 SAR therefore allows us to capture vessel activity completely absent from AIS. 
\end{quote}

\point{R3.9}
% category: writing/minor | assignee: Gavin
% Line 160: define "S1 footprint".
\begin{refquote}\itshape
Line 160: the term ``S1 footprint'' should be clearly defined. Does it refer to
the spatial coverage of a single S1 acquisition, the union of all acquisitions
over a certain period, or something else?
\end{refquote}
\response Thank you for this comment, and for the opportunity to clarify our methods. We now explicitly define the `S1 footprint' as the spatial union of all S1 scenes during our 2017 -- 2025 study period. We have also further clarified the text around the spatiotemporal hetreogeneity of S1 coverage and how we account for this explicitly in our modeling framework. For example, each pixel is imaged by S1 by varying amounts of observation effort each month due to differing rates of satellite overpass frequency and pixel coverage, so we create a normalized S1 observation effort metric that accounts for this to consistently measure vessel density across all pixel-months.

\changes The introduction now reads:

\begin{quote}
 \textcolor{red}{Lines 172-190}: To estimate unmatched vessel detection density over space and time, we must first account for the fact that S1 coverage is spatiotemporally heterogeneous. For pixels-months that are imaged at all by S1, we estimate vessel density by normalizing the number of vessels detected in that pixel-month by the total area that was imaged by S1 in the pixel month, with this area representing the sum across the per-scene areas of all spatial scene intersections with that pixel; this therefore accounts for satellite overpass frequency and coverage provides a normalized and consistent observation effort metric across all pixel-months. Meanwhile, the S1 footprint only ever covers certain regions of the globe, so there are some areas that are never imaged at all (here we define the `S1 footprint' as the spatial union of all S1 scenes during our 2017 -- 2025 study period). And even for those areas within the S1 footprint, different areas are imaged by S1 scenes at different frequencies over time, which means that not all areas within the S1 footprint are even imaged every month. To account for this heterogeniety, we build an additional set of machine learning models to predict unmatched vessel density for areas that are not imaged by S1 in a given month (either because this area is within the S1 footprint but did not receive S1 coverage in this particular month; or because this area falls completely outside the S1 footprint).
\end{quote}

\point{R3.10}
% category: writing/analysis | assignee: Zihan and David
% Figure 1: explain why S1-derived contribution is flat after 2020.
\begin{refquote}\itshape
Figure 1: The main text should explain why the S1-derived contribution remains
unchanged after 2020. Is this due to data-quality or data-availability issues,
or do the authors assume that emissions from non-broadcasting vessels no longer
increase and that all additional emissions come from broadcasting vessels? This
point is critical for interpreting the temporal evolution in the figure.
\end{refquote}
\response The `S1' trend shown in Fig. 1 represents the emissions from S1 vessel detections that could not be matched with AIS-broadcasting vessels. This trend is the joint result of two key trends: 1) the trend of total vessel activity (reflected by the trend of total emissions estimated by this study); and 2) the trend of detection-AIS match rate  (the fraction of S1 detections matched to AIS-broadcasting vessels). It is not the result of any assumptions, or  data-quality or data-availability issues. It should also be noted that S1 vessel detections unmatched with AIS could correspond to \textit{either} vessels that are completely non-broadcasting (i.e., they do not carry an AIS transponder at all); \textit{or} to vessels that carry AIS transponders but whose signal was not received for any number of reasons (i.e., broadcasting but signal-not-received). The `S1' line in Figure 1 therefore captures both of these types of vessels. The `S1' line excludes emissions from S1 vessel detections that were matched to AIS-broadcasting vessels; those emissions are already captured in the `AIS' line.

During the 2017-2025 study period, especially after 2020, global AIS adoption and reception quality has significantly increased, as shown by Global Fishing Watch's study: \url{https://globalfishingwatch.org/2025-one-ocean-science-congress-poster-rise-in-public-vessel-tracking/}. As a result, the fraction of S1 detections matched with AIS-broadcasting steadily increases after 2020 but shows much less change across 2017-2020 (see Fig. 11). From 2022 to 2023, the increase rate of matched fraction is 6.3\% and exceeds the increase rate of total vessel activities 5.0\%. As a result, the emission contributed by unmatched detections decreases from 2022 to 2023 while the total emission increases.

After 2020, although the AIS-derived emission steadily increases while S1-derived emission fluctuates, this does not mean that emissions from non-broadcasting vessels no longer increase and that all additional emissions come from broadcasting vessels. One of the factors driving increase in AIS activities (thus AIS-derived emissions) is the increasing AIS adoption, i.e., non-broadcasting vessels start broadcasting. However, another major factor is the improving AIS reception via both improvement in satellite reception quality and the advent of dynamic-AIS, which makes more activities by AIS-broadcasting vessels visible in AIS data. The latter scenario essentially moves part of the S1-derived emission to AIS-derived emission while not affecting the total estimated emission.
\changes We have clarified this in numerous places throughout the text. We clarified the caption for Figure 1 (Lines XX-YY) with the following text:

\begin{quote}
Annual marine vessel CO$_{2}$ emissions, with each line colored based on the data source: `AIS' (emissions from AIS-broadcasting vessels whose signal was received); `S1' (emissions from vessels detected by S1 that could not be matched to an AIS-broadcasting vessel, either because the vessel is completely non-broadcasting or because its AIS signal was not received); and `AIS + S1` (the aggregation of these `AIS' and `S1' emissions).
\end{quote}

We clarified the caption for Figure 2 (Lines XX-YY) with the following text:

\begin{quote}
(A) Time series of monthly global CO$_2$ emissions, with each line colored based on the data source: `AIS' (emissions from AIS-broadcasting vessels whose signal was received); `S1' (emissions from vessels detected by S1 that could not be matched to an AIS-broadcasting vessel, either because the vessel is completely non-broadcasting or because its AIS signal was not received); and `AIS + S1` (the aggregation of these `AIS' and `S1' emissions). (B) Time series of the monthly share of global CO$_{2}$ emissions attributable to S1-detected unmatched vessels. (C) Spatial distribution of cumulative 2017-2025 CO$_{2}$ emissions from AIS-broadcasting vessels (top), S1-detected unmatched vessels (middle), and both types of vessels combined (bottom), shown at a 1x1 degree spatial resolution; emissions are area-normalized for each pixel (metric tonnes per square kilometer, shown on a pseudolog-10 scale).
\end{quote}

\point{R3.11}
% category: figure | assignee: Gavin
% New/added figure: spatial distribution of S1-derived emissions.
% Cross-ref R2.6.
\begin{refquote}\itshape
The spatial distribution of ship emissions derived from S1 data is not presented.
Showing such a distribution, at least in aggregated form, would be very
informative and could be considered for inclusion (either in Figure 1 or as an
additional figure). This would also help readers understand where
non-broadcasting emissions are concentrated.
\end{refquote}
\response This is an excellent suggestion. We made a revised version of Figure 2 which now explicitly shows the spatial distribution of ship emissions derived from S1 data (middle map of panel C). To go along with this map, we also highlight in the text the areas where non-broadcasting emissions are concentrated. We also added a new Figure S1, which is similar to Figure 2, but faceted to show the differences between non-fishing and fishing vessels, the finest level of vessel class categorization detail possible with the S1 data.
\changes We revised Figure 2, and also updated the corresponding results paragraph text (Lines XX-YY). We also added a new Figure S1.

\point{R3.12}
% category: analysis/figure | assignee: Gavin (and anyone else with thoughts)
% Ship-type analysis: show added value of new dataset (distributions by type/
% region, broadcasting vs non-broadcasting). Cross-ref R3.13, R3.14.
\begin{refquote}\itshape
The analysis by ship type is rather confusing. Since the study introduces a new
dataset and methodology, readers would naturally expect to see how this new
information refines our understanding of ship-type composition and activity
patterns (e.g., distributions by ship type and region, differences between
broadcasting and non-broadcasting fleets). However, the manuscript focuses
instead on the composition of non-fishing and fishing vessels in two separate
figures, without clearly demonstrating the added value of the new dataset for
ship-type analysis.
\end{refquote}
\response We greatly appreciate your comment, and the opportunity to make our ship type analysis more clear and insightful. We have completely revised Figure 3. Figure 3A now shows ranked emissions by vessel class within four different families of vessels (cargo, tanker, service and passenger, and fishing), allowing the reader to quickly see differences between vessel classes of a particular family. These bar charts still have colors for international trips, domestic trips, and port visits, allowing the reader to see how emissions vary across these different activity types. Panel 3B is completely new, and shows the spatial distribution of emissions for each vessel class family, allowing the reader to quickly see differences in the regional patterns across these vessel types. We would also like to clarify that Figure 3 focuses only on AIS-broadcasting emissions, since we are only able to disaggregate vessel classes to this level of detail using our AIS data. Our non-broadcasting emissions data derived from S1 only allows us to disaggregate between non-fishing and fishing vessels, which as you rightly point out is not as valuable as having more refined vessel class information. We have clarified in the figure caption and throughout the figure panels themselves that these data focus only on AIS-based emissions.
\changes We have completely revised Figure 3, as well as the accompanying text in the results section (Lines XX-YY).

\point{R3.13}
% category: analysis/writing | assignee: Pol
% Figure 3a: compare ship-type contributions to other inventories (container vs
% bulk expectation); if driven by S1, analyse S1 ship-type composition. Ties R3.1.
\begin{refquote}\itshape
Figure 3a: previous studies generally indicate that container ships are the largest contributors to global ship CO\textsubscript{2} emissions, followed by bulk carriers. The authors should compare their ship-type contributions with other inventories and explain the reasons for any discrepancies. If the differences are primarily driven by the inclusion of S1 data, then the ship-type composition of S1-derived emissions should be analyzed and discussed more explicitly.
\end{refquote}
\response We thank the reviewer for this comment.

First, a clarification that resolves the second half of the comment: the differences in ship-type contributions are not driven by the inclusion of S1 data. Figure 3 reports AIS-based emissions only. This is a limitation of the S1 data rather than a choice. S1 detections support disaggregation only between fishing and non-fishing vessels, not into the 37 vessel classes of our AIS-based inventory, and we have now stated this explicitly in the figure caption and in the accompanying text, since the original figure did not make it clear. The ship-type composition in Figure 3 therefore reflects our AIS-based fleet alone, and any departure from other inventories must be explained by the AIS-based model, not by data fusion.

Second, on the comparison itself: within the cargo and tanker families our ranking agrees with previous studies. Container ships and bulk carriers are the highest-emitting cargo classes in our inventory, and oil tankers the highest-emitting tanker type, consistent with STEAM, MariTEAM, SEIM and SAVE. We now say so in the text, with citations, rather than reporting our ranking in isolation. The main discrepancy between our inventory and the other inventories is that, in 2025, passenger vessels rather than container ships constitute the largest single emitting group. We have two explanations for this, and both are now in the manuscript. The first concerns how we define this class in our model. `Passenger' is a highly aggregated category in our vessel classification, comprising ferries, cruise ships, Ro-Pax and yachts, which are treated as separate classes or excluded entirely in other inventories. This aggregation makes passenger the most numerous class in our fleet and consequently the largest emitting group. The second is fleet coverage, and it is the same mechanism that explains our divergence from other inventories in total emissions (R3.7): vessels below 25 m account for most passenger emissions in our inventory. These are small, largely coastal vessels that registries cover poorly, and registry-anchored inventories cannot represent. The same is true of fishing vessels, which are our second-largest class by vessel count and which most global shipping inventories exclude or underrepresent.


\changes We have revised the relevant Results paragraph, the caption of Figure 3, and added supporting supplementary material.

1. The Results paragraph on vessel classes now makes the comparison with other inventories explicit and explains the discrepancy.

2. The caption of Figure 3 now states that the figure reports AIS-broadcasting emissions only, and the panels themselves are labeled accordingly, so that no reader can attribute the vessel-class composition to the S1 component.

3. Two supplementary figures supporting the explanation: Supplementary Fig. 7, passenger vessels by length bin for 2017 against 2025, which shows the growth and the emissions concentrated in the smallest sizes; and Supplementary Fig. 8, fleet composition in 2025, giving each vessel class's share of the fleet against its share of CO$_{2}$, which shows directly that passenger's position is a function of vessel count.


\point{R3.14}
% category: writing/analysis | assignee: Gavin
% Fishing vessels international vs domestic framing: reconsider or justify.
\begin{refquote}\itshape
The discussion of fishing vessels in terms of international versus domestic trips
is not very meaningful. Fishing vessels are not transport vessels; their typical
activity pattern is to leave port, remain at sea for fishing operations, and then
return to port. Classifying their voyages as ``international'' or ``domestic''
trips is therefore of limited relevance and may be misleading. The authors should
reconsider whether this analysis is necessary and, if retained, justify its
relevance more clearly.
\end{refquote}
\response Thank you for this insightful comment. Since fishing vessels are not transport vessels, at your suggestion we have removed this description in the text. We therefore no longer describe how much fishing vessel emissions occurs during certain types of trips. We now simply say `Most emissions for cargo and tanker vessel classes occur during international trips.'
\changes

\point{R3.15}
% category: writing | assignee: Chris and Olivier
% Discussion reorg: move literature-review paragraphs to Intro; focus Discussion
% on implications/policy/limitations/future work. Coordinate with E1 (Discussion
% has no subheading).
\begin{refquote}\itshape
Discussion section: parts of the Discussion, especially the first two paragraphs,
read more like a literature review and would be more appropriate in the
Introduction. The Discussion should instead focus on the implications of the
results, including policy relevance, methodological limitations, and concrete
directions for future research. Reorganizing this section would improve the
clarity and impact of the manuscript.
\end{refquote}
\response We agree with this and have revised the entire discussion section along the lines of your suggestion.
\changes The discussion now focuses on policies that are supported by this work, externality costs that can be calculated, and temporal trends that can now be understood. We then turn to a discussion of limitations of the work and possibilities for overcoming them, including for future research.  

\end{document}
